\documentclass[10pt,letterpaper,sans]{moderncv}
\moderncvstyle{banking}
\moderncvcolor{blue}

\usepackage[scale=0.81]{geometry}
\nopagenumbers{}
\setlength{\hintscolumnwidth}{2.25cm}

\name{Krishnaa}{Vadivel}
\title{Software Engineer | Scientific Computing | HPC Workflows}
\email{srikrishnaa.careers@gmail.com}
\phone[mobile]{516-853-5663}
\social[linkedin]{sri-krishnaa-ece-phd}
\extrainfo{\href{https://github.com/krishnaa423}{GitHub: krishnaa423} \textbar{} \href{https://leetcode.com/u/krishnaa42342}{LeetCode: krishnaa42342} \textbar{} \href{https://hub.docker.com/u/krishnaa42342}{Docker Hub: krishnaa42342}}

\begin{document}
\makecvtitle

\section{Summary}
\cvitem{}{Software engineer and Yale PhD researcher building scientific-computing, automation, and high-performance workflows in Python and C++. Experience spans distributed HPC, CUDA/OpenGL tooling, containerized development, numerical software, and production support around both research and engineering systems.}

\section{Experience}
\cventry{2021--present}{PhD Research Assistant}{Yale University, Electrical Engineering}{}{}{\begin{itemize}%
\item Develop Python/C++ workflows for first-principles and many-body simulation, using MPI, OpenMP, PETSc, SLEPc, CMake, and Linux-based automation across leadership-class HPC environments.
\item Build reusable tooling around distributed compute, scientific pipelines, and ML-assisted workflows for materials and molecular problems.
\item Work across research code, documentation, and theory-to-implementation translation for scalable computational workflows.
\item Co-authored a 2025 \textit{Journal of the American Chemical Society} paper and presented APS talks in 2024, 2025, and 2026, demonstrating research output alongside production-style software work.
\end{itemize}}
\cventry{2019--2021}{Software and Embedded Systems Engineer}{Probe Technologies Inc. / Weatherford Wireline Services}{}{}{\begin{itemize}%
\item Built CUDA/OpenGL software for waveform processing, interactive visualization, and diagnostic tooling for large engineering datasets.
\item Developed ASP-style services and Microsoft SQL-backed workflows for internal engineering and field-data systems.
\item Contributed embedded-adjacent engineering support, giving strong systems-level context for software used alongside deployed hardware.
\end{itemize}}
\cventry{2018--2019}{Photonics Technical Writer}{FindLight}{}{}{\begin{itemize}%
\item Wrote technical content translating specialized optics and photonics topics into clear engineering-facing explanations.
\end{itemize}}

\section{Selected Projects}
\cvitem{DOLFINx Semiconductor Workflow}{Python finite-element and PETSc-based workflow for Poisson and drift-diffusion style simulation.}
\cvitem{Agentic Scientific Automation}{Retrieval-driven tooling that reads documentation and papers to generate structured computational-science workflow inputs.}
\cvitem{Distributed Scientific Tooling}{Reusable software workflows for simulation setup, engineering analysis, and container-oriented research infrastructure.}
\cvitem{Scientific Visualization Tooling}{CUDA/OpenGL and Python-based utilities for inspecting and communicating large numerical datasets.}

\section{Education}
\cventry{2021--present}{PhD, Electrical Engineering}{Yale University}{}{}{}
\cventry{2023}{MPhil, Electrical Engineering}{Yale University}{}{}{}
\cventry{2023}{MS, Electrical Engineering}{Yale University}{}{}{}
\cventry{2017--2018}{MEng, Electrical and Computer Engineering}{Cornell University}{}{}{}
\cventry{2013--2017}{BS, Electrical and Computer Engineering}{Cornell University}{}{}{}

\section{Technical Skills}
\cvitem{Languages}{Python, C, C++, JavaScript, Bash, SQL}
\cvitem{Scientific Computing}{MPI, OpenMP, PETSc, SLEPc, NumPy, pandas, SciPy, SymPy, matplotlib}
\cvitem{Systems}{Linux, CMake, Docker, Docker Compose, Kubernetes, gcloud}
\cvitem{Platforms}{CUDA, OpenGL, Frontier, Frontera, Perlmutter}
\cvitem{ML / Automation}{PyTorch, scikit-learn, agentic workflows, scientific automation}

\end{document}
